\chapter{Conclusions}
\setcounter{section}{0}

\section{General Conclusions}

This project set out to answer whether a Windows desktop can be controlled autonomously from a single spoken natural language instruction while keeping the full execution pipeline on the user's own device. The answer obtained through the implementation and evaluation of Autopilot is positive within the task scope defined for this work. The final system can transcribe a spoken request, extract the user's intention, generate a deterministic PDDL plan, ground that plan in the graphical interface through local vision-language models, execute the required actions, and verify the final state without additional human intervention.

The main contribution of the project is the integration of local neural models with automated planning in a practical desktop automation pipeline. The results show that neither a purely command-based approach nor a purely reactive LLM agent is sufficient to satisfy all the requirements addressed in this report. Command-based systems such as Windows Voice Access can be fast and reliable, but they require the user to decompose the task manually. Reactive cloud systems can accept open natural language, but they expose screenshots and task context to external services and may produce inconsistent action sequences. Fully local LLM-based systems preserve privacy, but the evaluation showed that tool access alone does not guarantee reliable automation. Autopilot demonstrates that a neuro-symbolic architecture can reduce these weaknesses by assigning each responsibility to the component best suited for it: language models interpret and orchestrate, vision-language models perceive the screen, and the symbolic planner generates the action structure.

The experimental evaluation supports this conclusion. Across 21 runs covering seven task templates, Autopilot achieved a 100\% task success rate, zero planning errors, zero tool hallucinations, zero data exposure, and a deterministic step structure. In the autonomous-system comparison, it matched Claude Cowork in task success while avoiding the data exposure of cloud-based systems and reducing planning errors. In the fully local comparison, it clearly outperformed OpenClaw, which did not complete any evaluated task under the same conditions. These results indicate that the architecture, rather than only the choice of models or tools, is the decisive factor behind the system's reliability.

At the same time, the project also identifies important limitations. The system is currently restricted to seven predefined PDDL task templates, so its determinism and reliability apply only inside that controlled task space. End-to-end latency remains high compared with command-based interaction, with an average of 111.83 seconds, mainly because local model inference and sequential orchestration are computationally expensive. Visual verification also remains imperfect: the QA Agent reported success in 18 of 21 runs, with the remaining errors caused by visual judgement limitations rather than actual task failure. These limitations do not invalidate the central hypothesis, but they show that local autonomous desktop automation is still a constrained and engineering-sensitive problem.

The qualitative study reinforces this balanced result. Participants perceived the system positively in terms of error recovery, behaviour, privacy, and accessibility, and all user-study tasks were completed successfully. However, daily-use likelihood and preference over alternatives were lower, mostly because of latency, transcription issues, and the limited number of supported templates. This suggests that Autopilot is already a strong proof of concept for private and accessible desktop automation, but not yet a general-purpose replacement for all forms of computer interaction.

From an ethical and social perspective, the project confirms the value of designing autonomous agents with privacy, auditability, and limited action scope from the beginning. Running all inference locally directly addresses data sovereignty concerns, while PDDL planning makes the intended action sequence inspectable before execution. The controlled template set also reduces the risk of arbitrary or harmful actions. The main ethical weakness identified is accountability, since the current system does not require explicit user confirmation before executing a plan. Overall, Autopilot shows that autonomy and privacy do not have to be treated as opposing goals, provided that the system architecture is designed around both from the start.

In conclusion, Autopilot validates the feasibility of a local, spoken, neuro-symbolic desktop automation system on consumer hardware. The project demonstrates that combining local LLMs, VLMs, speech recognition, deterministic planning, and GUI execution can produce a system that is autonomous, private, explainable, and empirically competitive within a defined task space. Its results are especially relevant for accessibility and privacy-sensitive contexts, where cloud-based automation or step-by-step command systems are not fully satisfactory.

\section{Conclusions regarding Project Objectives}

The main objective of the project was to develop a fully autonomous desktop automation system for Windows that operates from free-form spoken natural language, executes tasks through the graphical interface, verifies completion, and keeps all processing local. This objective was achieved within the seven supported task templates. The system completed all 21 evaluated runs successfully, worked from spoken input, required no further human intervention after the initial prompt, and transmitted no data outside the device.

\textbf{SO1: Data Sovereignty and Privacy.} This objective was fully achieved. The complete pipeline runs locally, including speech transcription, language understanding, visual perception, planning, orchestration, execution, and verification. The comparative evaluation assigns Autopilot a data exposure level of 0, in contrast with cloud-based alternatives such as Claude Cowork and Gemini Computer Use, both of which expose screenshots and task context externally. This confirms that the system satisfies the privacy-by-design requirement established at the beginning of the project.

\textbf{SO2: Robustness, Predictability, and Determinism in the Planning Layer.} This objective was also achieved. The PDDL planning layer produced deterministic plans with a planning error rate of 0 and a standard deviation of steps of 0. The use of Fast Downward and explicit task templates made the generated action sequences inspectable, reproducible, and free from missing or redundant planning actions in the evaluated scenarios. The main limitation is that this guarantee only applies to the predefined PDDL domains, not to arbitrary desktop tasks outside the supported task set.

\textbf{SO3: Closed-Loop Autonomy.} This objective was achieved for the evaluated scope. After the initial spoken command, Autopilot executed each task without additional keyboard, mouse, touch, or voice input from the evaluator. The system used Scout, Vision, Executor, Manager, and QA agents to observe the desktop, act on the GUI, recover from intermediate states, and verify the result. The acceptance criterion of at least 19 successful runs out of 21 was exceeded, with 21 successful task completions. All runs stayed below the 300-second threshold, although the average latency of 111.83 seconds shows that speed remains one of the clearest limitations of the current implementation.

\textbf{SO4: Benchmarking with Other Solutions.} This objective was achieved through a comparative evaluation against Windows Voice Access, OpenClaw, Claude Cowork, Gemini Computer Use, and Gemini Computer Use with external screenshot injection. The benchmark used the same seven task templates, three runs per template, and a common metric set covering success rate, latency, error rates, consistency, steps, and data exposure. Autopilot was the strongest system among fully autonomous and fully local alternatives, and it was the only evaluated system to combine 100\% task success, zero planning errors, deterministic step structure, and zero data exposure.

\textbf{SO5: Spoken Natural Language User Interface (Hands-Free Interaction).} This objective was achieved. Every evaluated Autopilot run began from a free-form spoken prompt, transcribed locally through Whisper and interpreted by the Intention Agent. The user did not need to learn a fixed command vocabulary or manually decompose the task into interface-level instructions. The qualitative study also confirmed that participants could initiate tasks naturally by voice, although some transcription errors were observed and should be considered a current usability limitation.

\textbf{SO6: Graphical User Interface for Desktop Automation.} This objective was only partially achieved. The implemented interface is useful as a demonstration layer because it allows the evaluator to trigger the system, observe the automation process, and show that the underlying pipeline works without requiring direct terminal interaction. However, it should not be considered a complete end-user application. The GUI was not designed or validated for daily use by non-technical users, and it lacks the polish, robustness, feedback mechanisms, and interaction design that would be required for a real accessibility-oriented product. Therefore, the objective was achieved as a functional prototype, but not as a fully finished user-facing interface.

Overall, five of the six specific objectives were satisfied within the defined scope of the project, while the graphical interface objective was only partially completed. The strongest achievements are privacy preservation, deterministic planning, autonomous task completion, and empirical competitiveness against other systems. The main remaining weaknesses are latency, limited task coverage, speech transcription robustness, the maturity of final-state visual verification, and the lack of a production-ready interface for real daily use.

\section{Next Steps}

The first next step is to fine-tune Gemma4 for the specific orchestration and intent-classification tasks required by Autopilot. The current system uses a general-purpose local model, which is flexible but introduces latency and occasional inefficiency in the Manager and Intention stages. A task-specific fine-tuned model could reduce execution time, make agent dispatching more precise, and improve consistency across repeated runs.

The second improvement is to integrate the Whisper speech-to-text output more directly with the Intention Agent. At the moment, speech transcription and intent extraction are treated as separate stages, which increases the computational cost of the pipeline. A tighter integration between the STT component and the intent-classification step could reduce duplicated processing and make the initial phase of the system faster.

The third next step is to evaluate the system with physically impaired users. The current qualitative study provides useful feedback from general participants, but the main social motivation of the project is accessibility for users with motor difficulties. Testing with this target group would make it possible to assess whether Autopilot addresses real accessibility needs, which interaction barriers remain, and which tasks would be most valuable in daily use.

The final next step is to design and implement a custom GUI for end users. The current interface demonstrates the technical feasibility of the system, but it is not a complete product interface. A future version should include a more accessible and polished design, clearer execution feedback, error messages, status indicators, configuration options, and interaction flows suitable for non-technical users who may rely on the system regularly.
